\documentclass[10pt,a4paper]{amsart}
\usepackage[margin=19mm]{geometry}
\usepackage{amsmath,amssymb,mathtools}
\usepackage[scr=rsfs,cal=euler]{mathalfa}
\usepackage{xfrac}
\usepackage[unicode,hidelinks,bookmarks=false]{hyperref}
\newcommand{\ie}{i.e.\ }
\newcommand{\cf}{cf.\ }
\newcommand{\resp}{respectively\ }
\newcommand{\Subsection}{\subsection}
\providecommand{\nobreakdash}{\nobreak\hbox{-}}
\setlength{\emergencystretch}{2em}
\multlinegap0pt
\usepackage{bm}
\usepackage{bbm}
\usepackage{dsfont}
\usepackage{xspace}
\usepackage{cancel}
\usepackage{mathtools}
\usepackage{xcolor}
\usepackage{amsthm}
\makeatletter
\@ifundefined{theorem}{\newtheorem{theorem}{Theorem}}{}
\@ifundefined{lemma}{\newtheorem{lemma}[theorem]{Lemma}}{}
\makeatother
\title[Priestley Representation of Distributive Precontact Lattices]{Priestley Representation of Distributive Precontact Lattices}
\author{Sergio A. Celani, Luciana Valenzuela}
\date{Source: arXiv:2606.25131v1 (CC BY 4.0)}
\begin{document}
\maketitle

\noindent\emph{Source and licence.} Priestley Representation of Distributive Precontact Lattices. Version arXiv:2606.25131v1 (CC BY 4.0), distributed under the Creative Commons Attribution 4.0 International licence, as recorded in the arXiv licence metadata for this exact version and shown on the abstract page https://arxiv.org/abs/2606.25131v1. Full rights notice: licenses/arxiv-2606.25131-CC-BY-4.0.txt.\par\smallskip

\noindent\textit{Editorial scope.} Counted unit: the Lemma whose source label is \texttt{None} in the article (source0.tex lines 1014--1028, source label \texttt{None}), delivered together with the article's own complete original proof of it (source0.tex lines 1029--1081, one \texttt{proof} environment of 53 lines). No mathematics is written, repaired or re-derived: statement and proof are the article's own text line for line. Required context retained uncounted: propp (lines 386--403). Every definition, display and label of those lines is retained unchanged. Omitted from this package and not redistributed: 1--385, 404--1013, 1082--2244. Nothing inside a retained excerpt is omitted or altered: every retained body is delivered exactly as the article's lines are written, so no omission receipt of any kind accompanies this package. Citations: the retained excerpts carry no citation command at all, so no bibliography entry of the article is transcribed and no reference is redistributed. Reference adaptation: none was needed and none was made; every reference inside the delivered unit resolves inside it, so no unresolved reference or citation is produced. Printed slips kept literal: no printed word, symbol, hypothesis or number of the counted statement or of its proof was altered. Layout and class: the article's own class is \texttt{article} with packages including fontenc, inputenc, color, colortbl, mathtools, enumitem, amsmath, amsthm, amssymb, hyperref, tikz, graphicx, array; the wrapper is the lane's pinned \texttt{amsart} editorial wrapper at 10pt on A4, which restates the theorem-like environments the retained text needs and adds the article's own macro definitions that the retained text needs (none), with extra wrapper packages (bm, bbm, dsfont, xspace, cancel, mathtools, xcolor). The delivered numbering is the unit's own. Assets: no retained span contains a figure environment, a graphics-inclusion command, a PDF-page inclusion, a table or a TikZ picture; no image file is copied into this package, the package declares no asset dependency and no image file is redistributed with it. Licence: version arXiv:2606.25131v1 of this article is distributed under the Creative Commons Attribution 4.0 International licence, as recorded in the arXiv licence metadata for this exact version and shown on the abstract page https://arxiv.org/abs/2606.25131v1. A full rights notice is delivered at licenses/arxiv-2606.25131-CC-BY-4.0.txt. Characters: every retained excerpt is pure ASCII, so the delivered engine, XeLaTeX with its default Latin Modern text font, reports no missing character.
\begin{lemma}\label{propp} Let $\langle L,\mathsf{C}\rangle$ be
a precontact lattice, and let $a,b\in L$. Then 
\begin{enumerate}
\item $\langle L,\mathsf{C}^{-1}\rangle$ is a precontact lattice. 
\item If $(a,b)\in\mathsf{C}$, $a\leq a'$ and $b\leq b'$ then $(a',b')\in\mathsf{C}$. 
\item If $a\neq0$ and $(a,a)\in\mathsf{C}$ then $(a,1)\in\mathsf{C}$. 
\item The sets $L\setminus\mathsf{C}(a)=\{b\in L:(a,b)\notin\mathsf{C}\}$
and $L\setminus\mathsf{C}^{-1}(a)=\{x\in L:(x,a)\notin\mathsf{C}\}$
are ideals of $L$. 
\item If $(a,b)\in\mathsf{C}$, then $[a)\times[b)\subseteq\mathsf{C}$. 
\item If $F$ and $G$ are filters of $L$ and $F\times G\subseteq\mathsf{C}$,
then there exist prime filters $P$ and $Q$ such that $F\subseteq P$,
$G\subseteq Q$, and $\textcolor{blue}{P\times Q \subseteq\mathsf{C}}$. 
\item For all $a,b\in L$, if $a\mathsf{C}b$, then there exist prime filters
$F$ and $G$ such that $F\times G\subseteq\mathsf{C}$, $a\in F$,
and $b\in G$. 
\end{enumerate}
\end{lemma} 

\begin{lemma} Let $\langle{L}_{1},\mathsf{C}_{1}\rangle,\langle{L}_{2},\mathsf{C}_{2}\rangle$
be a precontact lattice and let $h:{L}_{1}\longrightarrow{L}_{2}$
be a lattice homomorphism. Consider the map 
\[
h_{*}:{\rm {X}}(L_{2})\longrightarrow{\rm {X}}(L_{1})
\]
defined by $h_{*}(P)=h^{-1}(P)$, for each $P\in{\rm {X}}(L_{2})$.
Then $h_{*}$ is a Priestley morphism, i.e. continuos and order preserving,
such that: 
\begin{enumerate}
\item $h$ is a precontact homomorphism iff $h_{*}$ is a stable morphism. 
\item $h$ is a strong precontact homomorphism iff $h_{*}$ is a strong
stable morphism. 
\end{enumerate}
\end{lemma} 

\begin{proof}
$(1)$ Assume that $h$ satisfies condition $\mathrm{\mathrm{(CH1)}}$.
Let $(P,Q)\in R_{2}$. We aim to prove that $(h_{*}(P),h_{*}(Q))\in R_{1}$,
i.e. , $h^{-1}(P)\times h^{-1}(Q)\subseteq C_{1}$. Consider $(p,q)\in h^{-1}(P)\times h^{-1}(Q)$.
Then, $(h(p),h(q))\in P\times Q$. Since $P\times Q\subseteq\mathsf{C}_{2}$
by hypothesis, it follows that $(h(p),h(q))\in C_{2}$. By the stability
of $h$, we conclude that $(p,q)\in C_{2}$.

Assume that $h_{*}$ satisfies $\mathrm{(SP1)}$. Let $(h(a),h(b))\in\mathsf{C}_{2}$.
Then, by $(7)$ of Lemma \ref{propp}, there exists $P,Q\in{\rm {X}}(L_{2})$
such that $h(a)\in P,h(b)\in Q$ and $(P,Q)\in R_{2}$. As $(P,Q)\in R_{2}$
and $h_{*}$ is stable, $(h^{-1}(P),h^{-1}(Q))\in R_{1}$. Thus, $(a,b)\in h^{-1}(P)\times h^{-1}(Q)\subseteq C_{1}$.
So, $(a,b)\in C_{1}$.

$(2)$ We assume that $h\colon\langle L_{1},\mathsf{C}_{1}\rangle\longrightarrow\langle L_{2},\mathsf{C}_{2}\rangle$
satisfies conditions $\mathrm{\mathrm{(CH1)}}$ and $\mathrm{\mathrm{\mathrm{\mathrm{(CH2)}}}}$.
By $(1)$, $h_{*}$ satisfies $\mathrm{(SP1)}$. Therefore, it only
remains to prove that $h_{*}$ satisfies $\mathrm{(SP2)}$. Consider
$P\in\mathrm{X}(L_{2})$ and $Q\in\mathrm{X}(L_{1})$ such that $(h_{*}(P),Q)\in R_{1}$,
that is, $h^{-1}(P)\times Q\subseteq\mathsf{C}_{1}$.

We prove that $P\times\mathrm{Fg}(h(Q))\subseteq\mathsf{C}_{2}.$
Assume that there exist $p\in P$ and $x\in Q$ such that $(p,h(x))\notin\mathsf{C}_{2}$.
Then, by condition $\mathrm{(CH2)}$, there exists $b\in L_{1}$ such
that $(b,x)\notin\mathsf{C}_{1}$ and $p\leq h(b)$. However, this
leads to a contradiction, because $p\leq h(b)$ and $p\in P$ imply
that $b\in h^{-1}(P)$. Thus, $(b,x)\in h^{-1}(P)\times Q\subseteq\mathsf{C}_{1}$,
i.e. , $(b,x)\in\mathsf{C}_{1}$, which is a contradiction. Thus,
$P\times\mathrm{Fg}(h(Q))\subseteq\mathsf{C}_{2}.$ By item $(6)$
of Lemma \ref{propp}, there exist $D,F\in\mathrm{X}(L_{2})$ such
that $\mathrm{Fg}(h(Q))\subseteq D$, $P\subseteq F$, and $(F,D)\in R_{2}$.
This implies that $(P,D)\in R_{2}$. Therefore, there exists $D\in\mathrm{X}(L_{2})$
such that $(P,D)\in R_{2}$ and $Q\subseteq h^{-1}(D)$.

Now let us prove that if $h_{*}$ is a strong stable morphism, then
$h$ is a strong precontact homomorphism. By the previous theorem,
and since $h_{*}$ is a stable map, $h$ is a precontact homomorphism.
Therefore, it remains to show that $h$ satisfies $\mathrm{(CH2)}$.
Assume that $(h_{*}(P),Q)\in R_{1}$ implies that there exists $D\in{\mathrm{X}}(L_{2})$
such that $(P,D)\in R_{2}$ and $Q\subseteq h^{-1}(D)$. Let $(a,h(x))\notin\mathsf{C}_{2}$.
Suppose that for all $b\in L_{1}$, if $a\leq h(b)$ then $(b,x)\in\mathsf{C}_{1}$.
Hence, $h^{-1}(\mathrm{Fg}(a))\times\mathrm{Fg}(x)\subseteq\mathsf{C}_{1}$.
Thus, by $(6)$ of Lemma \ref{propp}, there exist $F,Q\in{\mathrm{X}}(L_{1})$
such that $h^{-1}(\mathrm{Fg}(a))\subseteq F$, $[x)\subseteq Q$,
and $(F,Q)\in R_{1}$. Since $h^{-1}(\mathrm{Fg}(a))\subseteq F$,
there exists $P\in{\mathrm{X}}(L_{1})$ such that $h^{-1}(P)\subseteq F$
and $\mathrm{Fg}(a)\subseteq P$. As $h^{-1}(P)\subseteq F$ and $(F,Q)\in R_{1}$,
it follows that $(h^{-1}(P),Q)\in R_{1}$. Then, applying the hypothesis,
there exists $D\in{\rm {X}}(L_{2})$ such that $(P,D)\in R_{2}$ and
$Q\subseteq h^{-1}(D)$. Consequently, $(a,h(x))\in\mathsf{C}_{2}$,
which is a contradiction. Therefore, there exists $b\in L_{1}$ such
that $(b,x)\notin\mathsf{C}_{1}$ and $a\leq h(b)$. 
\end{proof}

\end{document}
